\chapter{Całka Kurzweila--Henstocka}
\providecommand{\atlasref}[2]{#2}
\label{ma1:ch:kurzweil-henstock}\label{it:sec:kurzweil-henstock}
\begingroup
\providecommand{\osc}{\operatorname{osc}}
\providecommand{\Var}{\operatorname{Var}}
\providecommand{\mesh}{\operatorname{mesh}}
\providecommand{\supp}{\operatorname{supp}}
\providecommand{\AC}{\mathrm{AC}}
\providecommand{\BV}{\mathrm{BV}}
\providecommand{\HK}{\mathrm{KH}}
\providecommand{\ind}[1]{\chi_{#1}}
\providecommand{\Dn}{\underline D}
\providecommand{\Du}{\overline D}
\providecommand{\KS}{\mathrm{KS}}
Rozwijamy jednowymiarową teorię całek cechowanych, całek Perrona i Denjoya oraz całek Stieltjesa. Część elementarna wymaga jedynie całki Riemanna i rachunku różniczkowego. Od miejsca, w którym używamy miary Lebesgue'a, zakładamy znajomość podstaw teorii miary: regularności miary Lebesgue'a, twierdzeń granicznych dla całki Lebesgue'a, twierdzenia Radona--Nikodyma i różniczkowania całek oraz miar. Odpowiednie dowody znajdują się w skrypcie \emph{Teoria miary}, w szczególności w rozdziałach 07, 08, 16--18, 24 i 25. Symbol $\lambda_1$ oznacza miarę Lebesgue'a na prostej. Dla ograniczonej funkcji $F$ i niepustego podprzedziału $I$ piszemy $\osc_I F=\sup_I F-\inf_I F$. Gdy $F$ jest ciągła, a $I$ zwarty, supremum i infimum są osiągane, więc $\osc_I F=\max_I F-\min_I F$. Symbole $\AC$ i $\BV$ oznaczają klasy funkcji bezwzględnie ciągłych i o ograniczonym wahaniu. Dla dowolnej funkcji $F$ określonej na $I=[u,v]$ przyjmujemy
\[
 \Var_I F=\Var_u^v(F)=V_u^v(F)=\sup_{u=x_0<\cdots<x_m=v}\sum_{i=1}^m|F(x_i)-F(x_{i-1})|,
\]
z wartością zero dla $u=v$. Dla $P=\{x_0,\ldots,x_m\}$ oznaczamy przez $\mesh P=\max_{1\le i\le m}(x_i-x_{i-1})$ średnicę podziału. Dla ciągłej funkcji $F$ na przedziale zwartym jej moduł ciągłości to $\omega_F(h)=\sup_{|s-t|\le h}|F(s)-F(t)|$, gdzie $s,t$ należą do jej dziedziny. Funkcja $F$ jest bezwzględnie ciągła na $[a,b]$, jeżeli dla każdego $\varepsilon>0$ istnieje $\eta>0$ takie, że dla każdej skończonej rodziny przedziałów $[u_j,v_j]$ o rozłącznych wnętrzach z $\sum_j(v_j-u_j)<\eta$ zachodzi $\sum_j|F(v_j)-F(u_j)|<\varepsilon$. Przez $\mathcal R[a,b]$ oznaczamy funkcje całkowalne w sensie Riemanna, a przez $\mathcal L^1$ konkretne funkcje całkowalne w sensie Lebesgue'a; $L^1$ oznacza ich klasy równości prawie wszędzie.
 
\paragraph{Konwencje.}
O ile nie zaznaczamy inaczej, rozważamy funkcje rzeczywiste
o skończonych wartościach, niekoniecznie ograniczone, i przedział
zwarty $[a,b]$ z $a<b$. Dla funkcji zespolonej $f=u+iv$
całkowalność w danej teorii skalarnej oznacza całkowalność
obu części rzeczywistych, a całkę definiujemy jako
$\int u+i\int v$. W teorii Lebesgue'a dla funkcji nieujemnych
dopuszczamy również całkę nieskończoną.
Całkowalność zależy od funkcji, dziedziny, miary lub funkcji
całkującej oraz rodzaju całki. W szczególności przynależność
do klasy funkcji całkowalnych w sensie Riemanna nie jest
własnością klasy modulo równość prawie wszędzie.
Przez funkcję pierwotną w sensie KH rozumiemy funkcję
$F(x)=C+(\HK)\int_a^x f(t)\,\mathrm dt$, gdzie $C\in\R$;
ciągłość i równość $F'=f$ prawie wszędzie są twierdzeniami,
a nie częścią definicji.
Terminologia angielska: \emph{gauge} --- cechowanie,
\emph{tagged partition} --- podział znaczony,
\emph{refinement} --- rozdrobnienie podziału,
\emph{integrand} --- funkcja podcałkowa,
\emph{integrator} --- funkcja całkująca,
\emph{bounded variation} --- wahanie ograniczone,
\emph{regulated function} --- funkcja regulowana.

\begin{lemat}[Addytywność wahania i rozkład Jordana]\label{ma1:kh:lem:variation-additivity}\label{it:lem:rs-variation}
Jeśli $g\in\BV[a,b]$ i $a\le u\le c\le v\le b$, to
\[
 V_u^v(g)=V_u^c(g)+V_c^v(g).
\]
W szczególności $V(t)=V_a^t(g)$ jest niemalejąca, $V(v)-V(u)\ge |g(v)-g(u)|$, a funkcje $(V\pm(g-g(a)))/2$ są niemalejące i ich różnica jest równa $g-g(a)$.
\end{lemat}
\begin{proof}
Wstawienie $c$ do dowolnego podziału $[u,v]$ nie zmniejsza sumy modułów przyrostów, na mocy nierówności trójkąta. Po rozcięciu podziału w $c$ suma jest nie większa niż $V_u^c(g)+V_c^v(g)$. Biorąc supremum, otrzymujemy jedną nierówność. Dla $\varepsilon>0$ wybieramy podziały obu części, których sumy są odległe od odpowiednich supremów o mniej niż $\varepsilon/2$; ich połączenie daje sumę większą niż $V_u^c(g)+V_c^v(g)-\varepsilon$. Dowolność $\varepsilon$ daje równość. Przypadek przedziału jednopunktowego jest bezpośredni. Stosując równość do $[a,v]$ rozciętego w $u$, otrzymujemy $V(v)-V(u)=V_u^v(g)\ge |g(v)-g(u)|$ (używamy podziału złożonego z samych końców). Ta nierówność daje monotoniczność obu funkcji, a ich różnica wynika z definicji.
\end{proof}
\section{Cechowania i definicja całki}
\begin{definicja}
\label{ma1:kh:def:kh-gauge}\label{it:def:kh-gauge}
Dla $a<b$ podział znaczony to skończona rodzina komórek $[x_{i-1},x_i]$, gdzie $a=x_0<x_1<\cdots<x_m=b$, ze znacznikami $t_i\in[x_{i-1},x_i]$. Jego suma ma postać
\[
 S(f,\mathcal P)=\sum_{i=1}^m f(t_i)(x_i-x_{i-1}).
\]
Podział częściowy składa się ze skończenie wielu podprzedziałów $[a,b]$ o dodatnich długościach, parami rozłącznych wnętrzach i znacznikach należących do odpowiednich przedziałów. Na przedziale jednopunktowym przyjmujemy pusty podział znaczony, którego suma wynosi zero, i całkę równą zeru.


\emph{Cechowaniem} na $[a,b]$ nazywamy dowolną funkcję $\delta:[a,b]\to(0,\infty)$. Podział znaczony $\mathcal P=\{(t_i,[x_{i-1},x_i])\}$ jest $\delta$-drobny, jeśli $t_i\in[x_{i-1},x_i]\subset(t_i-\delta(t_i),t_i+\delta(t_i))$ dla każdego $i$.
\end{definicja}
\begin{lemat}[Cousina]\label{ma1:kh:lem:kh-cousin}\label{it:lem:kh-cousin}
Niech $a<b$. Dla każdego cechowania $\delta$ na $[a,b]$ istnieje $\delta$-drobny podział znaczony tego przedziału.
\end{lemat}
\begin{proof}
Przypuśćmy, że przedział nie ma takiego podziału. Dzielimy go na połowy; co najmniej jedna również nie ma drobnego podziału, bo dwa drobne podziały połówek można połączyć. Powtarzając, dostajemy zstępujący ciąg domkniętych przedziałów $J_n$ długości dążących do zera, z których żaden nie ma podziału drobnego. Z twierdzenia o zstępującym ciągu przedziałów domkniętych, których długości dążą do zera, ich przecięcie składa się z jednego punktu $t$. Dla dostatecznie dużego $n$ mamy $J_n\subset(t-\delta(t),t+\delta(t))$, więc pojedynczy znaczony przedział $(t,J_n)$ jest drobny. Sprzeczność.
\end{proof}
\begin{definicja}
\label{ma1:kh:def:kh-integral}\label{it:def:kh-integral}
Funkcja $f:[a,b]\to\R$ jest całkowalna w sensie Kurzweila--Henstocka z całką $I$, jeżeli dla każdego $\varepsilon>0$ istnieje cechowanie $\delta$ takie, że dla każdego $\delta$-drobnego podziału znaczonego $\mathcal P$ zachodzi $|S(f,\mathcal P)-I|<\varepsilon$.
\end{definicja}
\begin{stwierdzenie}[Jednoznaczność]
\label{ma1:kh:prop:kh-unique}\label{it:prop:kh-unique}
Wartość całki KH jest wyznaczona jednoznacznie.
\end{stwierdzenie}
\begin{proof}
Jeśli $I,J$ spełniają definicję, dla $\varepsilon>0$ bierzemy minimum cechowań odpowiadających błędom $\varepsilon/2$ dla obu wartości. Lemat Cousina zapewnia istnienie podziału drobnego dla tego minimum. Dla jego sumy $S$ mamy $|I-J|\le |I-S|+|S-J|<\varepsilon$. Dowolność $\varepsilon$ daje $I=J$.
\end{proof}
\begin{twierdzenie}[Kryterium Cauchy'ego]\label{ma1:kh:thm:kh-cauchy}\label{it:thm:kh-cauchy}
Funkcja $f$ jest całkowalna w sensie KH wtedy i tylko wtedy, gdy dla każdego $\varepsilon>0$ istnieje cechowanie $\delta$ takie, że $|S(f,\mathcal P)-S(f,\mathcal Q)|<\varepsilon$ dla wszystkich $\delta$-drobnych podziałów $\mathcal P,\mathcal Q$.
\end{twierdzenie}
\begin{proof}
Konieczność wynika z nierówności trójkąta. Dla dostateczności wybieramy cechowania $\delta_n$ odpowiadające błędom $2^{-n}$, zastępujemy je kolejno minimami z poprzednimi i wybieramy $\delta_n$-drobne podziały $\mathcal P_n$ z lematu Cousina. Dla $m\ge n$ mamy $\delta_m\le\delta_n$, więc oba podziały $\mathcal P_n,\mathcal P_m$ są $\delta_n$-drobne i $|S(f,\mathcal P_n)-S(f,\mathcal P_m)|<2^{-n}$. Sumy tworzą zatem ciąg Cauchy'ego. Niech $I$ będzie jego granicą. Dla danego $n$ i dowolnego $\delta_n$-drobnego podziału $\mathcal P$ porównujemy jego sumę z sumami $S(f,\mathcal P_m)$, $m\ge n$, i przechodzimy z $m\to\infty$. Otrzymujemy $|S(f,\mathcal P)-I|\le2^{-n}$. Dla dowolnego $\varepsilon>0$ wybieramy $n$ z $2^{-n}<\varepsilon$; cechowanie $\delta_n$ spełnia definicję całki.
\end{proof}
\begin{twierdzenie}[Liniowość, porządek i lokalność]\label{ma1:kh:thm:kh-locality}\label{it:thm:kh-locality}
Klasa funkcji KH-całkowalnych jest przestrzenią wektorową, całka jest liniowa i monotoniczna. Jeśli $a<c<b$, to $f$ jest całkowalna w sensie KH na $[a,b]$ wtedy i tylko wtedy, gdy jest całkowalna na obu $[a,c]$ i $[c,b]$; wówczas $\int_a^b f=\int_a^c f+\int_c^b f$.
\end{twierdzenie}
\begin{proof}
Dla funkcji całkowalnych $f,g$ i liczb $\alpha,\beta\in\R$ wybieramy minimum cechowań dla $f,g$ z błędem $\varepsilon/(1+|\alpha|+|\beta|)$. Liniowość sum daje
\[
 |S(\alpha f+\beta g,\mathcal P)-(\alpha I_f+\beta I_g)|
 <\varepsilon.
\]
Jeżeli $f\le g$, to dla każdego podziału $S(f,\mathcal P)\le S(g,\mathcal P)$. Wybierając wspólne cechowanie dla błędu $\varepsilon/2$ i podział z lematu Cousina, uzyskujemy $I_f\le I_g+\varepsilon$. Dowolność $\varepsilon$ daje monotoniczność.

Jeżeli $f$ jest całkowalna na $[a,b]$, ustalamy cechowanie $\delta$ z kryterium Cauchy'ego dla błędu $\varepsilon$. Na $[a,c]$ wybieramy jeden $\delta$-drobny podział $\mathcal R$ z lematu Cousina. Dla dwóch dowolnych drobnych podziałów $\mathcal P,\mathcal Q$ przedziału $[c,b]$ ich połączenia z $\mathcal R$ są drobnymi podziałami $[a,b]$. Po odjęciu wspólnego składnika $S(f,\mathcal R)$ otrzymujemy $|S(f,\mathcal P)-S(f,\mathcal Q)|<\varepsilon$. Kryterium Cauchy'ego daje całkowalność na $[c,b]$; zamiana ról obu przedziałów daje całkowalność na $[a,c]$. Stosując ten argument dwukrotnie, otrzymujemy całkowalność na każdym podprzedziale. W drugą stronę wybieramy cechowania na obu podprzedziałach dla błędu $\varepsilon/2$, dla $t\ne c$ dodatkowo mniejsze od $|t-c|$, a w punkcie $c$ bierzemy minimum cechowań. Każdy przedział drobnego podziału przecinający $c$ ma znacznik $c$; rozcinamy go w $c$. Powstałe części dają drobne podziały z tą samą sumą, ponieważ przy rozcięciu zachowujemy znacznik $c$, a długości dodają się. Każda z dwóch sum różni się od odpowiedniej całki lokalnej o mniej niż $\varepsilon/2$. Nierówność trójkąta daje błąd mniejszy od $\varepsilon$ względem sumy całek lokalnych, co dowodzi całkowalności i addytywności. Jednoznaczność zapewnia tę samą równość także wtedy, gdy zaczynaliśmy od całkowalności na całym przedziale.
\end{proof}
\begin{stwierdzenie}
$\mathcal R[a,b]\subset\HK[a,b]$ i obie całki są równe.
\end{stwierdzenie}
\begin{proof}
Z twierdzenia \ref{ma1:thm:riemann-darboux-equivalence} dla każdego $\varepsilon>0$ istnieje $\delta_0>0$ takie, że wszystkie sumy Riemanna dla podziałów o średnicy mniejszej od $\delta_0$ różnią się od całki Riemanna o mniej niż $\varepsilon$. Wybieramy cechowanie stałe $\delta(t)=\delta_0/2$. Każda komórka drobnego podziału ma długość mniejszą od $2\delta(t)=\delta_0$, więc jego suma spełnia żądane oszacowanie.
\end{proof}
\begin{przyklad}[Funkcja Dirichleta jest KH-całkowalna]\label{ma1:kh:ex:dirichlet}
Niech $\Q\cap[0,1]=\{q_1,q_2,\dots\}$ i $D=\chi_{\Q\cap[0,1]}$. Dla $\varepsilon>0$ kładziemy $\delta(q_k)=\varepsilon 2^{-k-3}$, a w punktach niewymiernych np. $\delta(t)=1$. Grupujemy wszystkie komórki o tym samym znaczniku $q_k$. Ich wnętrza są rozłączne, a wszystkie leżą w $(q_k-\delta(q_k),q_k+\delta(q_k))$, więc ich \emph{łączna} długość nie przekracza $2\delta(q_k)$, również gdy znacznik jest wspólnym końcem dwóch komórek. Stąd dla każdego drobnego podziału $0\le S(D,\mathcal P)\le\sum_k 2\delta(q_k)<\varepsilon$. Całka KH jest równa $0$.
\end{przyklad}
\begin{lemat}[Zmiany na zbiorze miary zero]\label{ma1:kh:lem:kh-null-modification}\label{it:lem:kh-null-modification}
Jeśli $h:[a,b]\to\R$ ma skończone wartości i jest równa zeru poza zbiorem $N$ miary Lebesgue'a zero, to jej całka KH istnieje i wynosi $0$. Zmiana funkcji na zbiorze miary zero nie zmienia jej całkowalności w sensie KH ani wartości całki, jeśli istnieje.
\end{lemat}
\begin{proof}
Rozbijamy $N$ na zbiory $N_k=\{t\in N:k-1<|h(t)|\le k\}$, $k\ge1$. Dla ustalonego $\varepsilon>0$ wybieramy otwarte pokrycia $U_k\supset N_k$ z $\lambda_1(U_k)<\varepsilon2^{-k-1}/k$. Jeśli $t\in N_k$, dobieramy $\delta(t)>0$ tak, aby $(t-\delta(t),t+\delta(t))\subset U_k$; w pozostałych punktach, również w punktach $N$ o $h(t)=0$, cechowanie jest dowolne, gdyż ich wkład do sumy wynosi zero. W podziale drobnym komórki z oznaczonymi punktami należącymi do $N_k$ mają rozłączne wnętrza i leżą w $U_k$. Ich wkład do modułu sumy jest co najwyżej $k\lambda_1(U_k)$. Sumując po $k$, dostajemy błąd mniejszy od $\varepsilon$.
 Jeśli dwie funkcje o skończonych wartościach różnią się tylko na zbiorze zerowym, ich różnica spełnia powyższy warunek. Liniowość całki KH (twierdzenie \ref{ma1:kh:thm:kh-locality}) daje ostatnią tezę, w obu kierunkach.
\end{proof}

\begin{center}
\begin{tikzpicture}[x=1.25cm,y=1cm,>=Stealth]
 \draw[->] (-.25,0)--(7.4,0);
 \draw[very thick,blue!65!black] (0,0)--(2,0);
 \draw[very thick,teal!70!black] (2,0)--(4.8,0);
 \draw[very thick,blue!65!black] (4.8,0)--(7,0);
 \foreach \x/\name in {0/{a=x_0},2/{x_1},4.8/{x_2},7/{x_3=b}}
   {\draw (\x,.10)--(\x,-.1) node[below] {$\name$};}
 \foreach \x/\name in {1/t_1,3.2/t_2,6/t_3}
   {\fill (\x,0) circle (2pt) node[above] {$\name$};}
 \draw[blue!65!black,<->] (-.2,.65)--(2.2,.65);
 
 \draw[teal!70!black,<->] (1.5,1.2)--(4.9,1.2);
 \draw[blue!65!black,<->] (4.6,.65)--(7.4,.65);
 \node[above] at (3.2,1.25) {$t_2\pm\delta(t_2)$};
\end{tikzpicture}

\small Rysunek: komórki podziału i otoczenia cechowania; ich rozmiary zależą od znacznika. Każda komórka musi zawierać się w swoim otwartym otoczeniu.
\end{center}
\section{Pochodne i lemat Saksa--Henstocka}
\begin{przyklad}[Wszystkie pochodne są całkowalne KH]\label{ma1:kh:ex:all-derivatives}
Jeśli $a<b$, a $F:[a,b]\to\R$ jest ciągła i różniczkowalna na $(a,b)$, to funkcja $f$ równa $F'$ we wnętrzu i mająca dowolne skończone wartości w $a,b$ jest KH-całkowalna oraz $\int_a^b f(x)\,\mathrm dx=F(b)-F(a)$. Dla $\varepsilon>0$ przy każdym $t\in(a,b)$ wybieramy $\delta(t)$ tak, aby $|F(y)-F(t)-f(t)(y-t)|\le\varepsilon|y-t|/(4(b-a))$ dla bliskich $y$. Na komórce $[u,v]$ ze znacznikiem $t$ odejmujemy przyrosty $F(v)-F(t)$ oraz $F(u)-F(t)$; nierówność trójkąta daje
\[
 |F(v)-F(u)-f(t)(v-u)|
 \le \frac{\varepsilon}{4(b-a)}(|v-t|+|u-t|)
 =\frac{\varepsilon}{4(b-a)}(v-u).
\]
Dla $t=a$ i $t=b$ korzystamy tylko z ciągłości $F$: dobieramy $\delta(t)>0$ tak, aby $|F(y)-F(t)|+|f(t)||y-t|<\varepsilon/4$ przy $|y-t|<\delta(t)$, $y\in[a,b]$. Komórka ze znacznikiem $a$ lub $b$ musi mieć ten punkt jako koniec; każdego z tych znaczników można użyć co najwyżej raz. Łączny błąd komórek ze znacznikami wewnętrznymi jest nie większy niż $\varepsilon/4$, a komórek ze znacznikami brzegowymi mniejszy niż $\varepsilon/2$. Suma przyrostów $F$ teleskopuje się do $F(b)-F(a)$, więc całkowity błąd jest mniejszy od $\varepsilon$. Nie potrzeba istnienia pochodnych jednostronnych na końcach.
\end{przyklad}
\begin{przyklad}[Pochodna nieabsolutnie całkowalna]\label{ma1:kh:ex:conditional-derivative}
Niech $F(x)=x^2\cos(\pi/x^2)$ dla $x\in(0,1]$ i $F(0)=0$. Wtedy $F'(0)=0$ oraz $F'(x)=2x\cos(\pi/x^2)+(2\pi/x)\sin(\pi/x^2)$. Pochodna jest całkowalna KH i $\int_0^1F'(x)\,\mathrm dx=F(1)-F(0)=-1$. Natomiast rozszerzona całka Lebesgue'a $\int_0^1|F'(x)|\,\mathrm d\lambda_1(x)=\infty$:  składnik $2x\cos(\pi/x^2)$ ma całkowalny moduł, ograniczony przez $2x$. Dla drugiego składnika podstawienie $u=\pi/x^2$ daje
\[
 \int_c^1\frac{2\pi}{x}|\sin(\pi/x^2)|\,\mathrm dx
 =\pi\int_\pi^{\pi/c^2}\frac{|\sin u|}{u}\,\mathrm du.
\]
Na każdym przedziale $[k\pi,(k+1)\pi]$, $k\ge1$, całka po prawej bez czynnika $\pi$ jest nie mniejsza niż $2/((k+1)\pi)$. Szereg harmoniczny jest rozbieżny. Z nierówności $|A+B|\ge|B|-|A|$ wynika rozbieżność całki modułu całej pochodnej.
\end{przyklad}
\begin{lemat}[Saksa--Henstocka: uzupełnianie podziału częściowego]\label{ma1:kh:lem:kh-saks-henstock}\label{it:lem:kh-saks-henstock}\label{lem:hk-uzupelnianie}\label{ma1:kh:lem:hk-uzupelnianie}
Niech $f\in\HK[a,b]$, $F(x)=\int_a^x f$, a $\delta$ będzie cechowaniem, dla którego wszystkie $\delta$-drobne sumy na $[a,b]$ różnią się od $F(b)-F(a)$ o mniej niż $\varepsilon$. Jeśli $\mathcal Q=\{(t_j,[u_j,v_j])\}_{j=1}^m$ jest dowolnym $\delta$-drobnym podziałem częściowym z rozłącznymi wnętrzami, to
\[
\left|\sum_{j=1}^m\bigl(f(t_j)(v_j-u_j)-F(v_j)+F(u_j)\bigr)\right|\le\varepsilon.
\]
Ponadto suma modułów wyrazów w nawiasach jest nie większa niż $2\varepsilon$.
\end{lemat}
\begin{proof}
Uporządkujmy przedziały częściowego podziału. Domknięcia luk o dodatniej długości między komórkami oraz przy końcach $[a,b]$ tworzą co najwyżej $m+1$ przedziałów $J_k$ o rozłącznych wnętrzach. Wraz z komórkami podziału częściowego pokrywają $[a,b]$; luki jednopunktowe pomijamy, ponieważ nie wnoszą przyrostów. Dla dowolnego $\eta>0$ lokalność całki KH i lemat Cousina pozwalają wybrać na każdym $J_k$ podział znaczony, drobny względem $\delta$ oraz względem cechowania dającego błąd $\eta/(m+1)$ dla lokalnej całki po $J_k$. Połączenie tych podziałów z $\mathcal Q$ daje pełny podział $\delta$-drobny. Przyrosty $F$ na wszystkich komórkach sumują się do $F(b)-F(a)$. Błąd na $J_k$ jest mniejszy od $\eta/(m+1)$, zatem suma błędów na $\mathcal Q$ ma moduł nie większy niż $\varepsilon+\eta$. Przechodzimy z $\eta\downarrow0$.

Ten sam argument działa dla każdego podzbioru komórek $\mathcal Q$. Stosując go osobno do komórek o błędzie dodatnim i ujemnym, widzimy, że suma błędów dodatnich oraz modułów błędów ujemnych nie przekracza osobno $\varepsilon$. Ich suma jest więc nie większa niż $2\varepsilon$.
\end{proof}
\begin{lemat}[Pięciokrotne powiększenie przedziałów]\label{ma1:kh:lem:fivefold-cover}
Niech $\mathcal I$ będzie rodziną domkniętych przedziałów o dodatnich długościach ograniczonych z góry przez $L<\infty$. Istnieje co najwyżej przeliczalna podrodzina $\mathcal J$ przedziałów o parami rozłącznych wnętrzach taka, że
\[
 \bigcup_{I\in\mathcal I} I\subseteq\bigcup_{J\in\mathcal J}5J,
\]
gdzie $5J$ jest domkniętym przedziałem o tym samym środku co $J$ i pięciokrotnie większej długości.
\end{lemat}
\begin{proof}
Dzielimy rodzinę na klasy $2^{-k-1}L<|I|\le 2^{-k}L$, $k=0,1,\ldots$. W każdej kolejnej klasie wybieramy maksymalną względem inkluzji podrodzinę o rozłącznych wnętrzach, rozłącznych także z wnętrzami wcześniej wybranych przedziałów. Taka rodzina istnieje na mocy lematu Zorna, ponieważ suma łańcucha rodzin dopuszczalnych jest dopuszczalna. Suma wszystkich wybranych rodzin jest przeliczalna lub skończona: każdemu wnętrzu przypisujemy pierwszą należącą do niego liczbę wymierną w ustalonej enumeracji $\Q$, a rozłączność daje różnowartościowość tego przyporządkowania.
Jeśli $I$ z klasy $k$ nie został wybrany, jego wnętrze przecina wnętrze pewnego wybranego $J$ z klasy o indeksie co najwyżej $k$ (w przeciwnym razie można byłoby go dołączyć na etapie $k$). Wtedy $|J|>|I|/2$. Dla środków $c_I,c_J$ mamy $|c_I-c_J|<(|I|+|J|)/2$. Dla $x\in I$ stąd
\[
 |x-c_J|\le |x-c_I|+|c_I-c_J|<|I|+|J|/2<5|J|/2.
\]
Zatem $I\subset 5J$. Przedział wybrany zawiera się oczywiście we własnym pięciokrotnym powiększeniu. To dowodzi tezy także dla rodziny pustej.
\end{proof}
\begin{twierdzenie}[Prawie wszędzie pochodna pierwotnej]\label{ma1:kh:thm:kh-primitive-derivative}\label{it:thm:kh-primitive-derivative}
Jeśli $f\in\HK[a,b]$, a $F(x)=\int_a^x f$, to $F$ jest ciągła, $F'(x)=f(x)$ prawie wszędzie oraz $f$ jest mierzalna w sensie Lebesgue'a.
\end{twierdzenie}
\begin{proof}
Ustalmy $t\in[a,b]$ i $\eta>0$. Wybieramy cechowanie z lematu Saksa--Henstocka dla błędu $\eta/2$ i bierzemy $x$ z $|x-t|<\delta(t)$. Dla $x=t$ oszacowanie przyrostu jest bezpośrednie; dla $x\ne t$ pojedyncza komórka o końcach $x,t$ ze znacznikiem $t$ jest drobnym podziałem częściowym. Oszacowanie modułu sumy błędów przez $\eta/2$ (z dowodu lematu) daje
$|F(x)-F(t)-f(t)(x-t)|\le\eta/2$, niezależnie od kolejności $x,t$. Jeśli ponadto $|x-t|<\eta/(2(1+|f(t)|))$, to $|F(x)-F(t)|<\eta$. Dowodzi to ciągłości, także jednostronnej na końcach. Z lematu Saksa--Henstocka otrzymujemy następujące oszacowanie: dla ustalonego $r>0$ zbiór punktów, w których istnieją dowolnie małe przedziały $[u,v]$ zawierające $x$ i spełniające $|F(v)-F(u)-f(x)(v-u)|>r(v-u)$, ma miarę zewnętrzną zero. Istotnie, ustalmy dowolne $\varepsilon>0$ i wybierzmy cechowanie $\delta$ zapewniające błąd mniejszy od $\varepsilon$ w definicji całki na $[a,b]$. Lemat \ref{ma1:kh:lem:kh-saks-henstock} ogranicza wówczas sumę modułów błędów każdego drobnego podziału częściowego przez $2\varepsilon$. Wybieramy tylko takie przedziały ze znacznikiem $x$, które leżą w $(x-\delta(x),x+\delta(x))$. Nadal istnieją dowolnie krótkie przedziały tego rodzaju przy każdym punkcie rozpatrywanego zbioru, więc tworzą drobne pokrycie Vitaliego. Z lematu \ref{ma1:kh:lem:fivefold-cover} wybieramy co najwyżej przeliczalną rodzinę przedziałów o rozłącznych wnętrzach, których pięciokrotne powiększenia pokrywają zbiór. Każdemu wybranemu przedziałowi przypisujemy jeden znacznik, dla którego zachodzi nierówność definiująca zbiór wyjątkowy. Dla każdej skończonej podrodziny lemat Saksa--Henstocka daje $r\sum_j(v_j-u_j)<2\varepsilon$. Biorąc supremum po skończonych podrodzinach, otrzymujemy $\sum_j(v_j-u_j)\le2\varepsilon/r$. Z przeliczalnej podaddytywności miary zewnętrznej i wzoru na długość pięciokrotnie powiększonego przedziału wynika, że miara zewnętrzna zbioru jest nie większa niż $10\varepsilon/r$, a $\varepsilon$ jest dowolne. Stosujemy ten argument do $r=1/n$, $n\ge1$, i wykluczamy przeliczalną sumę otrzymanych zbiorów zerowych oraz końce $a,b$. Dla każdego pozostałego $x$ i każdego $n$ istnieje promień, poniżej którego żaden przedział zawierający $x$ nie narusza oszacowania z $r=1/n$. W szczególności przedział o końcach $x,x+h$ daje
\[
 \left|\frac{F(x+h)-F(x)}h-f(x)\right|\le1/n
\]
dla dostatecznie małych niezerowych $h$ obu znaków. To dowodzi $F'=f$ prawie wszędzie. Przedłużamy $F$ przez stałe poza $[a,b]$. Ciągłe funkcje $n(F(x+1/n)-F(x))$ zbiegają do $f(x)$ prawie wszędzie; ich granica górna jest mierzalna, a po zastąpieniu nieskończonych wartości zerem otrzymujemy skończoną funkcję mierzalną równą $f$ prawie wszędzie. Zupełność miary Lebesgue'a daje mierzalność $f$.
\end{proof}

\begin{uwaga}[Dlaczego ograniczamy każdą sumę częściową]
Nie wystarcza samo oszacowanie modułu sumy wszystkich błędów: dodatnie i ujemne błędy mogą się znosić. Dopiero możliwość wybrania dowolnego podzbioru komórek w lemacie \ref{ma1:kh:lem:hk-uzupelnianie} pozwala otrzymać sumę modułów. Jest ona potrzebna w dowodzie całkowania przez części i przy własnościach funkcji pierwotnych.
\end{uwaga}

\section{Całki niewłaściwe i osobliwości}
\begin{twierdzenie}[Hake'a: związek z całkami niewłaściwymi]\label{ma1:kh:thm:kh-hake}\label{it:thm:kh-hake}
Niech $a<b$ i niech $f:[a,b]\to\R$ ma skończone wartości. Wówczas $f\in\HK[a,b]$ wtedy i tylko wtedy, gdy $f\in\HK[c,b]$ dla każdego $a<c<b$ oraz istnieje skończona granica $L=\lim_{c\downarrow a}\int_c^b f$. W tym przypadku $\int_a^b f=L$, niezależnie od wartości $f(a)$.
\end{twierdzenie}
\begin{proof}
Jeśli $f\in\HK[a,b]$, lokalność daje całkowalność na $[c,b]$ i $\int_c^b f=F(b)-F(c)$, gdzie $F(x)=\int_a^x f$. Ciągłość $F$ w $a$ i $F(a)=0$ dają granicę $F(b)$. Dla dostateczności definiujemy $F(a)=0$ i $F(x)=L-\int_x^b f$ dla $x>a$. Funkcja jest ciągła, a jej przyrosty na przedziałach oddalonych od $a$ są lokalnymi całkami. Wybieramy ciąg $b=c_0>c_1>\dots\downarrow a$. Na $J_n=[c_n,c_{n-1}]$ wybieramy cechowanie $\delta_n$ zapewniające błąd definicyjny mniejszy od $\varepsilon2^{-n-3}$. Lemat Saksa--Henstocka daje wówczas sumę modułów błędów na podziale częściowym nie większą niż $\varepsilon2^{-n-2}$. Dla $t\in(c_n,c_{n-1})$ bierzemy minimum $\delta_n(t)$ oraz połowy odległości od obu końców. W punkcie $c_n$, $n\ge1$, bierzemy minimum $\delta_n(c_n)$, $\delta_{n+1}(c_n)$ i połowy odległości od $c_{n-1},c_{n+1}$. W punkcie $b$ wystarczy minimum $\delta_1(b)$ i $(b-c_1)/2$. Komórka przecinająca $c_n$ musi mieć znacznik $c_n$ i nie może sięgać kolejnego punktu podziału; po rozcięciu obie części są drobne względem lokalnych cechowań. W $a$ wybieramy $\delta(a)>0$ z $|F(y)|+|f(a)|(y-a)<\varepsilon/2$ dla $a<y<a+\delta(a)$, a poza $a$ zmniejszamy cechowania tak, aby przedziały z innymi znacznikami nie zawierały $a$. Przedział początkowy dowolnego drobnego podziału ma znacznik $a$, pozostałe można pogrupować według $J_n$ (i ewentualnie rozciąć w $c_n$). Rozcięcie nie zmienia sumy Riemanna ani sumy przyrostów $F$. Na każdym $J_n$ otrzymujemy podział częściowy (pomijamy ewentualne komórki jednopunktowe po rozcięciu), więc suma modułów błędów nie przekracza $\varepsilon2^{-n-2}$. Komórka początkowa wnosi błąd mniejszy od $\varepsilon/2$. Ponieważ $\sum_{n\ge1}2^{-n-2}=1/4$, całkowity błąd jest mniejszy od $\varepsilon$. Suma przyrostów wynosi $F(b)-F(a)=L$.
\end{proof}
\begin{przyklad}[Szeregi i całkowalność]
Dla $I_n=(2^{-n},2^{-n+1}]$ i ciągu $(a_n)_{n\ge1}$ kładziemy $f(x)=2^n a_n$ na $I_n$, a $f(0)=0$. Wówczas
\[
f\in\HK[0,1]\iff\sum_{n\ge1}a_n\ \text{jest zbieżny},\qquad
f\in\mathcal L^1[0,1]\iff\sum_{n\ge1}|a_n|<\infty.
\]
Na każdym $[c,1]$, $c>0$, funkcja jest schodkowa z co najwyżej skończoną liczbą skoków. Dla $c\in I_n$ całka po $[c,1]$ jest równa $\sum_{k=1}^{n-1}a_k+2^na_n(2^{-n+1}-c)$, przy czym ostatni składnik ma moduł co najwyżej $|a_n|$. Gdy szereg jest zbieżny, $a_n\to0$, więc ta całka ma granicę równą sumie szeregu. Odwrotnie, biorąc $c=2^{-n}$ otrzymujemy sumy częściowe szeregu. Twierdzenie Hake'a daje pierwszą równoważność; druga z $\int_{I_n}|f|\,\mathrm d\lambda_1=|a_n|$. Ponadto $f$ jest całkowalna w sensie Riemanna dokładnie wtedy, gdy $\sup_n2^n|a_n|<\infty$: konieczność wynika z ograniczoności każdej funkcji całkowalnej w sensie Riemanna, a dostateczność z ograniczoności i kryterium Lebesgue'a, ponieważ nieciągłości mogą wystąpić jedynie w przeliczalnym zbiorze $\{0\}\cup\{2^{-n}:n\ge0\}$. Dla $a_n=(-1)^{n+1}/n$ otrzymujemy $\int_0^1f=\log2$ w sensie KH, lecz $\int_0^1|f|\,\mathrm d\lambda_1=\infty$.
\end{przyklad}
\begin{przyklad}[Nie każda osobliwość jest całkowalna]
Funkcja $f(x)=1/x$ na $(0,1]$, niezależnie od skończonej wartości w $0$, nie jest całkowalna KH na $[0,1]$, bo $\int_c^1\mathrm dx/x=-\log c\to\infty$. Na $[-1,1]$ funkcja $1/x$ z dowolną skończoną wartością w $0$ ma wartość główną Cauchy'ego równą $0$, ale nie ma całki KH, gdyż jej ograniczenie do $[0,1]$ nie jest KH-całkowalne.
\end{przyklad}
\begin{twierdzenie}[Rodzina całek oscylacyjnych]\label{ma1:kh:tw:hk-oscylacyjne}\label{tw:hk-oscylacyjne}
Niech $p\in\R$, $q>0$ i niech $f_{p,q}:[0,1]\to\R$ będzie określona wzorem $f_{p,q}(0)=0$ oraz $f_{p,q}(x)=x^{-p}\sin(x^{-q})$ dla $x\in(0,1]$. Wówczas
\[
 f_{p,q}\in\HK[0,1]\quad\Longleftrightarrow\quad p<q+1,
 \qquad f_{p,q}\in\mathcal L^1[0,1]\quad\Longleftrightarrow\quad p<1.
\]
W szczególności dla $1\le p<q+1$ całka KH istnieje, lecz nie jest bezwzględnie zbieżna.
\end{twierdzenie}
\begin{proof}
Na $[\varepsilon,1]$ funkcja jest ciągła. Podstawiając $t=x^{-q}$, uzyskujemy
\[
\int_\varepsilon^1x^{-p}\sin(x^{-q})\,\mathrm dx
=\frac1q\int_1^{\varepsilon^{-q}}t^{(p-1)/q-1}\sin t\,\mathrm dt.
\]
Oznaczmy wykładnik przez $\beta=(p-1)/q-1$. Gdy $\beta<0$, funkcja $t^\beta$ maleje do zera. Dla $1\le A<B$ całkowanie przez części daje
\[
\int_A^B t^\beta\sin t\,\mathrm dt
=[-t^\beta\cos t]_A^B+\beta\int_A^B t^{\beta-1}\cos t\,\mathrm dt.
\]
Moduł prawej strony jest nie większy niż $A^\beta+B^\beta+|\beta|\int_A^B t^{\beta-1}\,\mathrm dt=2A^\beta$. Zatem kryterium Cauchy'ego daje zbieżność całki niewłaściwej. Gdy $\beta=0$, pierwotna $-\cos t$ nie ma granicy. Jeśli $\beta>0$, na półokresie $[k\pi,(k+1)\pi]$, $k\ge1$, sinus ma stały znak. Stąd
\[
 \left|\int_{k\pi}^{(k+1)\pi}t^\beta\sin t\,\mathrm dt\right|
 \ge 2(k\pi)^\beta,
\]
gdyż całka z $|\sin t|$ po tym półokresie wynosi $2$. Narusza to kryterium Cauchy'ego całki niewłaściwej. Twierdzenie Hake'a daje pierwszą równoważność, ponieważ $\beta<0$ dokładnie wtedy, gdy $p<q+1$.

Dla $\beta<-1$ całka modułu jest skończona przez oszacowanie $t^\beta|\sin t|\le t^\beta$ dla $t\ge1$. Dla $-1\le\beta<0$ mamy na każdym półokresie
\[
 \int_{k\pi}^{(k+1)\pi}t^\beta|\sin t|\,\mathrm dt
 \ge 2((k+1)\pi)^\beta.
\]
Szereg tych dolnych ograniczeń jest rozbieżny, ponieważ $(k+1)^\beta\ge 1/(k+1)$. Dla $\beta\ge0$ dolne ograniczenie wynosi $2(k\pi)^\beta\ge2$ i również daje rozbieżność. Zatem bezwzględna całkowalność zachodzi dokładnie dla $\beta<-1$, czyli $p<1$.
\end{proof}
\begin{przyklad}[Dokładny próg przy całce $x^{-2}\sin(x^{-2})$]
Dla $p=q=2$ mamy $2<3$, więc całka $\int_0^1x^{-2}\sin(x^{-2})\,\mathrm dx$ istnieje w sensie KH, ale nie w sensie Lebesgue'a. Podstawienie z dowodu daje użyteczną reprezentację obliczeniową
$\int_0^1x^{-2}\sin(x^{-2})\,\mathrm dx=\tfrac12\int_1^\infty t^{-1/2}\sin t\,\mathrm dt$.
\end{przyklad}
\endgroup
